\section{Introduction}
\label{sec:intro}

\IEEEPARstart{M}{obile} applications run neural networks on images, speech, and sensor data, but a full model can exceed the compute and memory budget of a phone. Split inference divides the model between the device and a nearby edge server. The device runs the first layers and sends an intermediate feature, and the edge runs the rest of the model~\cite{kang2017neurosurgeon,laskaridis2020spinn}. Split learning trains such a model across many users without collecting their raw data~\cite{gupta2018distributed,vepakomma2018split,thapa2022splitfed}. Personalized split learning adds a classifier at the end of the device part~\cite{han2023splitgp,rizwan2026splitomc}. Each request therefore has two classifiers available. The \emph{device exit} is personalized to the classes that its user collects, and the \emph{edge exit} is trained across all users of a cell.

Which of the two classifiers gives the right answer depends on where the user is. At home, most requests come from the classes that the user's own data contains. When the user commutes to another area, the device starts to receive requests from classes that are common there but absent from its own training data. In the commute scenario that we study, about one in nine requests of a device at home comes from such classes, and about one in two after the user leaves home. The two exits are accurate on opposite parts of this traffic. The device exit is correct on 86\% of requests from the user's own classes but on only 7\% of requests from classes learned by other users of the current cell. The edge exit is correct on 72\% and 41\% of the same requests. Neither exit is the better choice for the whole day.

A natural answer is to choose an exit for each request. Common split inference answers on the device when the device exit is confident and otherwise sends the request to the edge~\cite{han2023splitgp,laskaridis2020spinn,teerapittayanon2016branchynet}. This rule needs the device exit to be uncertain on classes it has not learned, and a personalized classifier does not behave that way. It places most of its probability on the user's own classes even for a request from another class, so its entropy separates the two kinds of requests only weakly (AUROC 0.64). As a result, confidence-based offloading sends 92\% of the requests to the edge and is no more accurate than always using the edge exit. It is even 1.3 points below always using the device exit.

Another answer is to change how strongly each device personalizes during training, for example by keeping more of the local model at home and more of the shared model while away~\cite{deng2020apfl}. We tested the best version of this idea, an oracle that knows whether each device is at home. The oracle raises the accuracy of devices at home by 3.7 points but lowers that of devices away from home by 1.7 points. Devices that keep more of their local model also make the cell's shared model more specialized to their own classes, and visitors who depend on that shared model lose accuracy. The gain of one device is paid for by others, so training-time adaptation gives little when users move independently.

These observations suggest that a device does not need to choose between its two exits. If the device averages the class probabilities of both exits, each exit can decide the answer on the requests where it is confident and correct. On a request from the user's own classes, the device exit puts a large probability on the right class. On a request from another user's class, the device exit spreads its probability over the user's classes, and the edge exit's confident answer remains the largest. A probability average keeps this property, while a product of probabilities does not, because a near-zero probability from the device exit removes the edge exit's correct answer. The average still has a problem. The edge exit is trained on the class mix of the whole cell, in which the device's own classes make up about one quarter of the samples. A device at home, however, receives these classes in about 88\% of its requests. The edge exit therefore undervalues exactly the classes that dominate the device's traffic.

We present DriftGate, an inference method for personalized split learning under user mobility. DriftGate answers with a weighted average of the class probabilities of the two exits. Before averaging, it removes the bias that the cell's class mix places on the edge exit, so that the device's own classes are no longer discounted. It sets the weight of each exit from how confident the two exits have been on the device's recent requests, which adapts the method to the model and the split point without labels. Finally, it sends only the requests on which the device exit is uncertain to the edge, so the device can trade accuracy for communication. DriftGate does not change training and adds a few microseconds of computation per request.

We evaluate DriftGate in eight settings. They include a commute scenario, real mobility traces from the GeoLife dataset~\cite{zheng2010geolife}, faster movement, partial participation, cell-wide changes in the class mix, random mobility, 100 classes, and a ResNet-18 split~\cite{he2016resnet}. We compare DriftGate with nine inference rules that cover confidence-based offloading, each single exit, output combination as in FedRoD~\cite{chen2022fedrod}, per-request weighting as in FedTHE~\cite{jiang2023fedthe}, early-exit ensembling as in ZTW~\cite{wolczyk2021ztw}, and label-shift correction~\cite{saerens2002adjusting}. DriftGate is the most accurate rule in all eight settings. It is 0.1 to 1.9 points above the strongest alternative in each setting and 5.4 to 12.5 points above confidence-based offloading. Devices in their home cell gain most, the lead over the better single exit is largest during the commute, and the gain extends to the devices with the lowest accuracy. In seven settings, DriftGate exceeds every alternative that sends all requests to the edge while sending at most half of them. \todo{Summary of on-device latency and energy on the smartphones and Jetson boards.}

This paper makes the following contributions.
\begin{itemize}
    \item We show how user mobility breaks inference in personalized split learning. The device exit and the edge exit are accurate on opposite requests, but choosing an exit by confidence cannot separate these requests, and adapting personalization during training moves accuracy from visitors to residents.
    \item We design DriftGate, which averages the probabilities of the two exits after removing the class bias of the edge exit, sets the weight of each exit without labels, and supports selective offloading. It runs on top of the trained model and does not change training.
    \item We evaluate DriftGate against nine inference rules in eight settings, including real mobility traces, and measure its cost on smartphones and embedded boards.
\end{itemize}

The rest of the paper is organized as follows. Section~\ref{sec:motivation} measures how the two exits and the two common remedies behave under mobility. Section~\ref{sec:related} reviews related work. Section~\ref{sec:system} defines the system and the inference problem, and Section~\ref{sec:design} presents DriftGate. Section~\ref{sec:evaluation} reports the evaluation, Section~\ref{sec:discussion} discusses the scope of the method, and Section~\ref{sec:conclusion} concludes.
